% ===== 原文 PDF pp.89--92（印刷 pp.88--91）=====
% tr-08 / 块 tag 5.33--5.36 / 文件 b
% 起点：p.89 第二句（p.88→p.89 的跨页句已由 tr08a 写完整）
% 终点：p.92 末（含跨到 p.93 的句子，已按「开始者写完」补完）
% 本文件含：表 6（p.89）、表 7（p.92）

% [原文 p.88]
由于7-膜在垂直于 D3 膜的方向上非紧，7-7 弦被冻结，于是可以只关注 4d 极限下的 3-3 弦与 3-7 弦：前者实现规范玻色子，后者实现规范理论中的荷电物质，而这些物质还带有来自7-膜的整体对称性 $F$。探针 D3 膜的数目对应于其世界体积上实现的 4d 有效规范理论的规范群 $G$ 的秩。当把 D3 膜放在表 6 所列 K3 面上某个 Kodaira 奇点的正上方时，D3 膜上的 4d 有效规范理论就到达一个共形点，于是出现一个味对称性为 $F$ 的 4d $\cN=2$ SCFT。事实上，理论的共形性也体现在：在这些奇点附近轴子--伸缩子 $\tau$ 是常数 [176]。

\begin{center}
\small
\begin{tabular}{lccccccc}
\toprule
Kodaira 型 & $II$ & $III$ & $IV$ & $I_0^*$ & $IV^*$ & $III^*$ & $II^*$ \\
\midrule
$F$ & $\varnothing$ & $\SU(2)$ & $\SU(3)$ & $\SO(8)$ & $E_6$ & $E_7$ & $E_8$ \\
$\Delta_7$ & $6/5$ & $4/3$ & $3/2$ & $2$ & $3$ & $4$ & $6$ \\
$\tau$ & $e^{\pi i/3}$ & $e^{\pi i/2}$ & $e^{\pi i/3}$ & free & $e^{\pi i/3}$ & $e^{\pi i/2}$ & $e^{\pi i/3}$ \\
\bottomrule
\end{tabular}
\end{center}
\tabtitle{6}{轴子--伸缩子 $\tau$ 取常数值时的各类 Kodaira 奇点，以及7-膜上相应的规范群 $F$。注意当 $F=\SO(8)$ 时，$\tau$ 是一个未被固定的常数，此时的背景不过是微扰 IIB 型定向折叠的 $\Z_2$ 对合，即 4 个 D7 膜叠放在一个 O7 平面之上。为了便于与后文的讨论衔接，对每个7-膜我们还列出了相应的 $\Delta_7$，它决定亏损角 $\frac{2\pi}{\Delta_7}$——该角量度绕7-膜一周时横向几何角坐标的改变。}

% [Tutor Note: 表 6 的列顺序是 Kodaira 纤维 II → III → IV → $I_0^*$ → $IV^*$ → $III^*$ → $II^*$，对应的规范群依次为 无 → SU(2) → SU(3) → SO(8) → E6 → E7 → E8，与 $\Delta_7=6/5,4/3,3/2,2,3,4,6$ 一一对应（$\Delta_7$ 与群的大小同向：群越大，亏损角 $2\pi/\Delta_7$ 越小）。$I_0^*$ 一行 $\tau$ 写作 free，是因为那里是微扰的 O7 平面加 4 个 D7 膜，$\tau$ 不被固定；其余各行 $\tau$ 取固定值，正是「共形点」的几何标志。]
% [Tutor Note: $\tau$ 列中 $e^{\pi i/2}$ 就是 $i$、$e^{\pi i/3}$ 则是 $\rho=e^{2\pi i/3}$ 的二分之一次方；这些值都是 $\SL(2,\Z)$ 的有限阶固定点，因此才允许「S 对偶 + 轨道折叠」不破坏共形性。后面表 7 会看到 $k=4$ 一行直接写 $i$，与这里的 $e^{\pi i/2}$ 是同一个数。]

让我们以单个 D3 膜在 F 理论中探测表 6 所列的 $I_0^*$ 奇点为例（推广到其他类型的奇点是直接的）。这样的设置实现了 4d $\cN=2$ 的 $\mathrm{U}(1)$ 规范理论，其八个荷电超多重态来自 D3 与 4 个 D7 膜之间的弦，以及后者在定向折叠作用下的镜像 [174,175,282]（更系统的描述见 [283]）。当这个 D3 膜位于 $I_0^*$ 奇点正上方时，环绕奇点的 3-3 弦带来的 $\Z_2$ 单值性把规范对称性 $\mathrm{U}(1)$ 提升为 $\SU(2)$，而味对称性 $F$ 也（从 D7 膜上的 $\mathrm{U}(4)$）提升为 $\SO(8)$。\footnote{这一提升也可以这样看出：把带一个 O7 平面、且在四个不动点各处放 4 个 D7 膜的 $T^2/\Z_2$ 上的 IIB 型理论 T-对偶到 $T^2$ 上的 I 型弦论。$T^2$ 上相应的 Wilson 线把 I 型的规范群 $\SO(32)$ 破缺为 $\SO(8)^4$，而每个 $\SO(8)$ 因子被解释为 IIB 型一侧 $T^2/\Z_2$ 的某个不动点处增强的规范对称性。进一步地，在此 T-对偶下，IIB 型的 D3 膜被映射为 I 型背景中卷绕 $T^2$ 的 D5 膜，从而由于填满时空的定向折叠平面而承载增强的 $\mathrm{USp}(2)$ 规范对称性。}所得的 SCFT 就是
% [原文 p.89]
带四个基本超多重态的 $\cN=2$ $\SU(2)$ 共形 SQCD。把 D3 膜沿垂直于7-膜的横向空间（即 K3 的底 $B_1:=\mathbb{P}^1$）移离 $I_0^*$ 奇点，我们便探测到该 SCFT 的库仑分支；确实，D3 膜上的世界体积理论约化为一个 $\mathrm{U}(1)$ 规范理论，而 $\tau$ 在 $B_1$ 上的全纯变化 [283] 与由 Seiberg--Witten 几何得到的库仑分支上有效规范耦合的行为定量吻合 [21]。此外，4d 场论中在 EFT 里扮演重要角色的规范瞬子，被等同于 IIB 型弦论中的 D$(-1)$ 瞬子；其非微扰效应自动编码在 K3 椭圆纤维化的轮廓之中 [283]。换句话说，人们可以把这一个秩一 4d SCFT 的 SW 曲线直接等同于 K3 在 $I_0^*$ 奇点被解消之后的椭圆纤维。另一方面，希格斯分支对应于把 D3 膜溶入 D7 膜上的规范磁通，这与一个 $\SO(8)$ 规范瞬子的约化模空间相符。

推广到秩 $r$ 也是直接的：只需考虑一叠 $r$ 个 D3 膜去探测上述7-膜（等价地，Kodaira 奇点）之一。此时 $r$ 维的库仑分支就简单地是对应秩一情形的“种子”库仑分支的对称轨道折叠（symmetric orbifold），\footnote{不过，更高秩理论的 SW 几何不能等同于 K3 的椭圆纤维化。}而希格斯分支则由 $r$ 个 $F$-瞬子的约化模空间给出，其复维数为 $2(rh_F^\vee-1)$，其中 $h_F^\vee$ 表示 $F$ 的对偶 Coxeter 数。

% [Tutor Note: $h_F^\vee$ 是对偶 Coxeter 数（$A_n:n+1$，$D_n:2n-2$，$E_6:12$，$E_7:18$，$E_8:30$）。以 $F=\SO(8)$（即 $D_4$，$h^\vee=6$）为例，公式给 $r=1$ 时复维数 $2(1\cdot 6-1)=10$，即四元数维 $5$；这与直接数一致：$\SU(2)$ 配 4 个基本超多重态（等价地 8 个半超多重态），其希格斯分支四元数维 $=4\times 2-\dim\SU(2)=8-3=5$。两条路给出同一个数，可用来记住公式里为什么出现 $h_F^\vee$。]

注意，上述被 D3 膜探测的 Kodaira 奇点全都是 crepant 的。于是自然要问：4d $\cN=2$ SCFT 能否由 F 理论中 D3 膜探测更奇异（但被允许）的奇点、例如终端奇点来构造？确实，近来 [57,58] 给出了在存在 D7 膜时借助所谓 S-折来构造 4d $\cN=2$ SCFT 的新方案，随后的 [59,284,285] 则论证它可以重现 [118,163,165,286] 中纯场论地分类出的全部秩一 4d SCFT。

稍后，[59,284] 又讨论了带离散挠（discrete torsion）的 S-折的推广，它生成的 4d SCFT 也可以由 [57,58] 中无挠的 S-折经希格斯化得到。总而言之，这类新方法为 4d SCFT 提供了大批构造，并且仍在研究之中，以期对 4d $\cN=2$ SCFT 的分类（至少是低秩的情形）带来新的启示。下面我们沿用 [57] 的思路，简要总结这类构造背后的想法。

% [原文 p.90]
F 理论中的 S-折\footnote{这些理论具有 $\SU(3)_r\times \mathrm{U}(1)_r$ 的 R 对称性，且其真空模空间上的无质量场与 $\cN=4$ 理论中的一致。}最初由 [263,287] 引入以研究 4d $\cN=3$ SCFT，它是 IIB 型弦论中定向折叠的一种（非微扰）推广：其几何上的轨道折叠作用伴随着作用在轴子--伸缩子 $\tau=C_0+i e^{-\phi}$ 上的非平凡 S 对偶作用，从而把 $\tau$ 固定到某个特定值。粗略地说，它等价于在横向几何 $(\C^3\times T^2)/\Z_k$ 上作 F 理论，其中 $k=2,3,4,6$。\footnote{立刻可以看出，$k=2$ 的 S-折退化为 IIB 型定向折叠紧化中通常的 $\mathrm{O3}^-$ 平面，后者有微扰描述。}

% [Tutor Note: 原文 $\tau=C_0+i e^{-\phi}$ 中 $i$ 与 $e$ 均为数学斜体（600 dpi 裁图目视核对，字形倾斜、非正体），与全书其余公式排法一致；勿误植为 $\mathrm{e}^{\mathrm{i}}$。]
为了看清 S-折如何工程化一个 4d $\cN=3$ 理论，我们首先看 $\C^3$，并用 $z_i$（$i=1,2,3$）记 $\C^3$ 的坐标，用 $z_4\equiv x+\tau y$ 记复结构为 $\tau$ 的 $T^2$ 上的坐标。$\Z_k$ 的作用则定义为 [57]
\begin{equation}
z_i \to e^{i\Phi_i} z_i\,,\qquad e^{i\Phi_i}\in\Z_k\,,
\tag{5.33}
\end{equation}
这预期会把探测 6d 横向空间 $\C^3$ 的 D3 膜所携带的 16 个超荷减去四分之一。\footnote{注意，F 理论中的纤维 $T^2$ 并不是物理时空的真正组成部分，因此不被 D3 膜探测。}要看出这一点，注意 $\C^3$ 的等距群给出 D3 膜世界体积上 4d $\cN=4$ 理论的 R 对称性 $\SU(4)_r$，16 个超荷 $Q_i$（$i=1,2,3,4$）在该群下按四维基本表示变换。于是，作用在 $\C^3$ 上的商 $\Z_k$ 自然诱导出超荷 $Q_i$ 上的一个 R 对称转动，它由下式生成 [57,288]
\begin{equation}
  \begin{aligned}
    r_k:\ & Q_i \to M^i{}_j Q_j\,,\\
    & M^i{}_j = \diag\Big(e^{i(\Phi_1+\Phi_2+\Phi_3)/2},\ e^{i(\Phi_1-\Phi_2-\Phi_3)/2},\ e^{i(-\Phi_1+\Phi_2-\Phi_3)/2},\ e^{i(-\Phi_1-\Phi_2+\Phi_3)/2}\Big)\,.
  \end{aligned}
  \tag{5.34}
\end{equation}

% [Tutor Note: (5.34) 里四个指数的符号组合是 $(+,+,+),(+,-,-),(-,+,-),(-,-,+)$——负号个数均为偶数。这正是 $\SO(6)_r\cong\SU(4)_r$ 的旋量表示 $\mathbf{4}$ 的四个权 $(\pm\Phi_1\pm\Phi_2\pm\Phi_3)/2$（取偶数个负号）；写成这种「半和」形式是因为超荷在 $\SU(4)_r$ 下是旋量，$\Z_k$ 作用在 $\C^3$ 上时只以一半的角度转动旋量。]
此外，$\Z_k$ 在 $T^2$ 上的作用也对超荷有某些影响。要成为 $T^2$ 的等距，群 $\Z_k$ 必须取 $k=2,3,4,6$。因此，除 $k=2$ 的情形外，$T^2$ 的复结构 $\tau$ 被限制到表 7 所列的值。由于 $\tau$ 取特殊值，$\SL(2,\Z)$ 的 S 对偶元素在超荷 $Q_i$ 上诱导出一个额外相位，在此设置下它作用为 [99,263,287,289]
\begin{equation}
s_k:\ Q_i \to e^{-\pi i/k} Q_i\,.
\tag{5.35}
\end{equation}

现在，若在 (5.34) 中选取合适的相位使得 $\Phi_1=\Phi_2=-\Phi_3=\frac{2\pi}{k}$，则 16 个超荷中有 12 个被由 $r_k s_k$ 生成的商 $\Z_k$ 保留下来，从而在探针 D3 膜上得到 4d $\cN=3$ 的 SCFT [263,287]。注意，所有这些由 S-折生成的（刚性）$\cN=3$ 理论都是非拉格朗日量的，因而内禀地强耦合。\footnote{注意，一个具有拉格朗日量描述的微扰 4d $\cN=3$ 理论必然提升为 $\cN=4$，例如参见 [290]。}关于这些 $\cN=3$ SCFT 的更多细节，参见 [263,287,291--294]。

% [Tutor Note: 把 (5.34) 与 (5.35) 相乘可验算：$\Phi_1=\Phi_2=-\Phi_3=2\pi/k$ 时，(5.34) 的四个相位为 $e^{i\pi/k},e^{i\pi/k},e^{i\pi/k},e^{-3i\pi/k}$，再乘上 (5.35) 的 $e^{-\pi i/k}$ 得 $1,1,1,e^{-4\pi i/k}$。于是当 $k=3,4,6$ 时最后一个相位 $\neq1$，恰好投掉 4 个超荷（$16\to12$，即 $\cN=4\to\cN=3$）；而 $k=2$ 时 $e^{-4\pi i/k}=1$，四个相位全为 1——这与脚注所说「$k=2$ 的 S-折退化为微扰的 $\mathrm{O3}^-$ 平面」是自洽的（微扰定向折叠不投掉超荷）。因此「$16\to12$」这句话实际是对 $k=3,4,6$ 而言。]

% [原文 p.91]
\begin{center}
\begin{tabular}{lcccc}
\toprule
 & $k=2$ & $k=3$ & $k=4$ & $k=6$ \\
\midrule
$\tau$ & free & $\tau=e^{\pi i/3}$ & $i$ & $\tau=e^{\pi i/3}$ \\[6pt]
$\rho$ & $\begin{pmatrix} -1 & 0 \\ 0 & -1 \end{pmatrix}$
       & $\begin{pmatrix} -1 & -1 \\ 1 & 0 \end{pmatrix}$
       & $\begin{pmatrix} 0 & -1 \\ 1 & 0 \end{pmatrix}$
       & $\begin{pmatrix} 0 & -1 \\ 1 & 1 \end{pmatrix}$ \\
\bottomrule
\end{tabular}
\end{center}
\tabtitle{7}{$\Z_k$ 商（$k=2,3,4,6$）所固定的轴子--伸缩子 $\tau$ 的取值。相应的 $\SL(2,\Z)$ 单值性矩阵 $\rho$ 也随之固定。}

% [Tutor Note: 表 7 的 $\rho$ 可逐格用「迹 / 行列式」核对：$k=2$ 是 $-I$（迹 $-2$，行列式 $1$，二阶）；$k=3$ 的两个行向量为 $(-1,-1)$ 与 $(1,0)$（迹 $-1$，三阶）；$k=4$ 为 $(0,-1)$ 与 $(1,0)$（迹 $0$，即 $S$，四阶）；$k=6$ 为 $(0,-1)$ 与 $(1,1)$（迹 $1$，六阶）。这四个矩阵在 $\SL(2,\Z)$ 中的阶恰为 $2,3,4,6=k$，说明 $\rho$ 正是把 $\tau$ 固定的那个元素（$\rho^k=1$ 是商能闭合的必要条件），也与 $k\Delta_7\in\{2,3,4,6\}$ 的约束同源。另注意 $k=4$ 一行的 $\tau$ 直接写作 $i$，即表 6 中的 $e^{\pi i/2}$。]

读者可能会发现，由 S 对偶扭转与 R 对称性 $\SO(6)_r$ 在超荷上诱导的两个变换 $s_k$ 与 $r_k$，一般而言并不会完全抵消；因此取不同的相位 $\Phi_i$，就有可能让探针 D3 膜上保留更少的超对称。例如，取
\begin{equation}
\Phi_1 = \frac{2\pi}{k}\,,\qquad \Phi_2 = -\Phi_3 \quad \mathrm{mod}\ 2\pi\Z\,,
\tag{5.36}
\end{equation}
则 D3 膜堆的世界体积保留 $\cN=2$ 超对称。从现在起，我们把这一类 S-折称为 $\cN=2$ S-折，以便与上面的 $\cN=3$ S-折相区分。

[57,58] 的要点在于：人们可以进一步把 $\cN=2$ S-折与 7-膜结合起来，生成更多的 4d $\cN=2$ SCFT。确实，当把 7-膜放在 $\C^3$ 中 $z_1=0$ 的轨迹上时，7-膜所保留的超荷，是在复平面 $z_1$ 上转动 $\SO(2)$ 的生成元下本征值为 $+1$ 的那些超荷——这与 S-折 (5.36) 的作用相容。然而这过于草率了，因为我们还没有考虑 7-膜对几何的反作用。一叠 7-膜会在平坦的横向平面中造成一个 $\frac{2\pi}{\Delta_7}$ 的亏损角，其中 $\Delta_7$ 与 $\tau$ 按表 6 相应地固定。我们用 $u$ 记这个锥奇点上的平坦坐标，它与插入 7-膜之前的 $\C$ 坐标 $z_1$ 的关系为 $z_1=u^{\Delta_7}$。于是 $z_1$ 的 $\Z_k$ 转动诱导 $u$ 转过角度 $\frac{2\pi}{k\Delta_7}$，从而诱导超荷转过一个相位 $e^{-\frac{\pi i}{k\Delta_7}}$。\footnote{单项式 $z_1^k$ 是 $\Z_k$ 不变的，它被等同于 D3 膜上 $\cN=2$ 理论中一个库仑分支（CB）算子的真空期望值，其标度维数为 $k\Delta_7$，这也正是它的 $\mathrm{U}(1)_R$ 荷。由此我们还能看出，$z_1$ 的 $\Z_k$ 转动在超荷上诱导出 $\mathrm{U}(1)_R$ 相位 $\omega$，它必须满足 $\omega^{k\Delta_7}=1$。}正如已经提到的，要保持超对称，就应当用一个 $\Z_{k\Delta_7}\subset\SL(2,\Z)$ 变换来补偿这个相位。于是必须有 $k\Delta_7=2,3,4,6$，从而得到以下六种可能 [57,58]：
% [边界] 本句在原文中跨 p.92→p.93（「…offset this phase by a transformation Z_{kΔ7} ⊂ SL(2,Z).」），已在此写完整；下一文件 tr08c 从 p.93 的 bullet 列表续写。

% [Tutor Note: 为什么是 $k\Delta_7=2,3,4,6$：$\SL(2,\Z)$ 的有限阶元素只有 $1,2,3,4,6$ 阶（阶 $n$ 的元素共轭于 $ST$ 的幂、$\pm I$、$S$ 等），要用 $\Z_{k\Delta_7}\subset\SL(2,\Z)$ 的元素去「抵消」$e^{-\pi i/(k\Delta_7)}$，其阶必须落在 $\{2,3,4,6\}$ 内、且不能是 1（阶 1 即无作用）。这条约束后面直接给出三种 $k$ 与 $\Delta_7$ 的配对，是本节最关键的算术。]
% [Tutor Note: 脚注里「is identified the VEV」原文漏了 with（600 dpi 裁图逐词核对，确为 "is identified the VEV"）；译文按「被等同于……的 VEV」处理。另注意这里的 $z_1^k$ 与前面的 $z_1=u^{\Delta_7}$ 一致：$z_1^k=u^{k\Delta_7}$ 是锥奇点上唯一 $\Z_k$ 不变的单值单项式，故它才是良定义的 CB 算子。]
